% Monte-Carlo-simulations.tex
% Topic: Monte Carlo Simulations
% (Advanced topic — outside the HSC Mathematics Extension 2 syllabus)

\textbf{Monte Carlo simulation} is a computational technique that uses random sampling to approximate numerical results. It is particularly useful when a problem is too complex to solve analytically.

\textbf{Basic idea:}
\begin{enumerate}[leftmargin=*]
  \item Define a random experiment whose outcome is related to the quantity you want to estimate.
  \item Simulate the experiment a large number of times (e.g.\ $N = 10{,}000$ trials).
  \item Estimate the quantity as the proportion (or average) of the simulated outcomes.
\end{enumerate}

\textbf{Estimating $\pi$ by Monte Carlo.} Consider a unit square with an inscribed quarter-circle of radius $1$.
\begin{itemize}[leftmargin=*]
  \item Generate $N$ random points $(x_i, y_i)$ with $x_i, y_i \in [0,1]$.
  \item Count the number $M$ of points satisfying $x_i^2 + y_i^2 \leq 1$.
  \item Then $\pi \approx 4M/N$.
\end{itemize}

\textbf{Why it works.} The probability that a random point in the unit square lies inside the quarter-circle is $\pi/4$, since the area of the quarter-circle is $\pi/4$ and the area of the square is $1$.

% Placeholder: further examples and simulations to be added.
